\documentclass[11pt,a4paper]{article}

% ── packages ──────────────────────────────────────────────────
\usepackage[T1]{fontenc}
\usepackage[utf8]{inputenc}
\usepackage[margin=2.5cm]{geometry}
\usepackage{graphicx}
\usepackage{booktabs}
\usepackage{amsmath,amssymb}
\usepackage{hyperref}
\usepackage{caption}
\usepackage{subcaption}
\usepackage{xcolor}

\hypersetup{colorlinks=true, linkcolor=blue, citecolor=blue, urlcolor=blue}

% ── figure path ───────────────────────────────────────────────
\graphicspath{{../checkpoints/curriculum/}}

% ── document ──────────────────────────────────────────────────
\begin{document}

% ============================================================
\section{Preliminary Experiment}
\label{sec:prelim}
% ============================================================

To validate the core game dynamics and verify that curriculum training can
induce a strategically coherent equilibrium, we conducted a preliminary
experiment using a lighter architecture---a multi-layer perceptron (MLP)
with REINFORCE---on a version of OracleGambit without the balance system.
Training proceeded in three phases (Fig.~\ref{fig:overview}).

\paragraph{Phase A -- Honest Host, Players Only.}
The host was fixed as a perfectly honest oracle ($\text{Honesty}=1.0$)
and only the players were trained.
Across 18{,}048 rounds, the follow rate rose from $0.25$ (random chance)
to $0.69$, with the win rate matching it exactly---confirming that players
successfully learned to exploit a truthful signal.

\paragraph{Phase B -- Deceptive Host, Host Only.}
Players were frozen and only the host was trained to maximise its own
reward by deceiving players.
Over rounds 20{,}096--42{,}112, signal honesty collapsed from $0.24$ to
$0.04$, while the host's win rate dropped from $0.24$ to $0.11$.
This seemingly counterintuitive result arises because the frozen players
continued to follow a signal that was now systematically false, driving
both parties toward anti-coordinated outcomes.

\paragraph{Phase C -- Joint Training (500 checkpoints).}
Both agents were unfrozen and trained simultaneously for over one million
rounds (44{,}160--1{,}042{,}176).
As shown in Fig.~\ref{fig:honesty_follow}, an \emph{adversarial phase}
dominated early Phase C: the host briefly exploited residual player trust
(follow rate $\approx 0.56$ at the phase start) before players adapted.
The gap $\Delta = \text{Follow Rate} - \text{Signal Honesty}$ oscillated
around zero and steadily contracted (Fig.~\ref{fig:delta}),
ultimately converging to the theoretical mixed-strategy Nash equilibrium
at $\text{Honesty} \approx \text{Follow Rate} \approx 0.25$---equal to
the random-chance baseline of $1/D$ ($D=4$ doors).

These results confirm that (i) the zero-sum payout formula is correctly
implemented, (ii) curriculum shaping is necessary to avoid trivial
convergence, and (iii) the system reliably finds the mixed-strategy
equilibrium under joint training.

% ---- Figure 1: Full Overview --------------------------------
\begin{figure}[htbp]
  \centering
  \includegraphics[width=\linewidth]{fig_overview.pdf}
  \caption{%
    \textbf{Curriculum training overview (MLP+REINFORCE, pre-balance).}
    Win rate, signal honesty, and follow rate across all three phases.
    Shaded regions indicate Phase~A (blue), Phase~B (orange), and
    Phase~C (green); dashed vertical lines mark phase transitions.
    The grey dashed line marks the random-chance baseline ($0.25$).
  }
  \label{fig:overview}
\end{figure}

% ---- Figure 2: Phase C Honesty vs Follow Rate ---------------
\begin{figure}[htbp]
  \centering
  \includegraphics[width=\linewidth]{fig_honesty_follow.pdf}
  \caption{%
    \textbf{Phase C -- Signal Honesty vs.\ Follow Rate.}
    Thin lines show raw per-checkpoint values; thick lines show
    MA-15 smoothed trends.
    Red shading indicates periods when players over-trust the host
    (\emph{follow $>$ honesty}); blue shading indicates under-trust.
    Both metrics converge to the mixed Nash equilibrium at $\approx 0.25$.
  }
  \label{fig:honesty_follow}
\end{figure}

% ---- Figure 3: Phase C Delta --------------------------------
\begin{figure}[htbp]
  \centering
  \includegraphics[width=\linewidth]{fig_delta.pdf}
  \caption{%
    \textbf{Phase C -- $\Delta = \text{Follow Rate} - \text{Signal Honesty}$.}
    The dotted line shows the MA-50 trend.
    The orange band marks the early adversarial phase (large oscillations);
    the blue band marks convergence toward Nash equilibrium ($\Delta \to 0$).
  }
  \label{fig:delta}
\end{figure}

\end{document}
